\documentclass[10pt,a4paper]{amsart}
\usepackage[margin=19mm]{geometry}
\usepackage{amsmath,amssymb,mathtools}
\usepackage[scr=rsfs,cal=euler]{mathalfa}
\usepackage{xfrac}
\usepackage[unicode,hidelinks,bookmarks=false]{hyperref}
\newcommand{\ie}{i.e.\ }
\newcommand{\cf}{cf.\ }
\newcommand{\resp}{respectively\ }
\newcommand{\Subsection}{\subsection}
\providecommand{\nobreakdash}{\nobreak\hbox{-}}
\setlength{\emergencystretch}{2em}
\multlinegap0pt
\usepackage{amssymb}
\usepackage[shortlabels]{enumitem}
\newtheorem{theorem}{Theorem}
\title{Spherical Milnor Spaces II: Projective Quotients and Higher Topological Structures}
\author{Jean-Pierre Magnot}
\date{Source: arXiv:2605.16370v1 (CC BY 4.0)}
\begin{document}
\maketitle

\noindent\emph{Source and licence.} Spherical Milnor Spaces II: Projective Quotients and Higher Topological Structures. Version arXiv:2605.16370v1 (CC BY 4.0), distributed by arXiv under the Creative Commons Attribution 4.0 International licence, http://creativecommons.org/licenses/by/4.0/, as recorded in the arXiv licence metadata for this exact version and shown on the abstract page https://arxiv.org/abs/2605.16370v1. Full licence notice: \texttt{licenses/arxiv-2605.16370-CC-BY-4.0.txt}.\par\smallskip

\noindent\textit{Editorial scope.} Counted unit: the article's theorem theorem (source0.tex lines 1010--1028). The article's complete original proof of it is delivered with it (source0.tex lines 1030--1145, one \texttt{proof} environment of 116 lines). No mathematics is written, repaired or re-derived: the statement and the proof are the article's own lines, in the article's own notation, and no retained line is substituted or re-flowed.  No further source-local context is needed: every cross-reference of every retained line resolves inside the two delivered spans.  Nothing was omitted from any retained span, so the package carries no omission record.  No retained line contains a citation, so the wrapper carries no bibliography and no bibliography entry.  No retained line contains a figure environment, an image-inclusion command or drawing code, so no image is delivered and the package declares no asset dependency.  Editorial formatting only, in addition: an AMS \texttt{amsart} preamble that restates the article's own declarations and macros used by the retained lines, namely the environments \texttt{theorem} on separate counters, together with the standard packages those lines need.  The retained environments print in this document as Theorem 1 in order of appearance, whereas the article prints the same items with the numbering of its own full document.  The complete native source of the article as distributed in the arXiv e-print for this exact version is bundled unchanged as \texttt{sources/200715/source0.tex}.
\begin{theorem}
The quotient map
\[
\pi:
E^{\mathrm{sph}}(G)
\longrightarrow
E^{\mathrm{proj}}(G)
=
E^{\mathrm{sph}}(G)/\mathbb Z_2
\]
defines a principal \(\mathbb Z_2\)-bundle in the diffeological sense.

Consequently, it determines a canonical characteristic class
\[
\alpha_{\mathrm{sph}}
\in
H^1(E^{\mathrm{proj}}(G);\mathbb Z_2).
\]
\end{theorem}

\begin{proof}
Recall that \(\mathbb Z_2=\{\pm1\}\) acts on
\(E^{\mathrm{sph}}(G)\) by
\[
\varepsilon\cdot(x_i,g_i)
=
(\varepsilon x_i,g_i).
\]

We first show that this action is smooth, free, and proper in the
diffeological sense.

\medskip

\noindent
\textbf{Smoothness.}
On each generating chart
\[
S_I\times G^I,
\]
the action is induced by
\[
\varepsilon\cdot((x_i)_{i\in I},(g_i)_{i\in I})
=
((\varepsilon x_i)_{i\in I},(g_i)_{i\in I}),
\]
which is smooth since multiplication by \(\pm1\) is smooth on the Euclidean
space \(\mathbb R^I\). By the definition of the final diffeology, the induced
action on \(E^{\mathrm{sph}}(G)\) is smooth.

\medskip

\noindent
\textbf{Freeness.}
Assume that
\[
\varepsilon\cdot(x_i,g_i)=(x_i,g_i)
\]
in \(E^{\mathrm{sph}}(G)\). By definition of the quotient relation, this means
that for every index \(i\) such that \(x_i\neq0\),
\[
\varepsilon x_i=x_i.
\]
Since
\[
\sum_i x_i^2=1,
\]
at least one coordinate satisfies \(x_i\neq0\). Hence
\[
\varepsilon=1.
\]
Therefore the action is free.

\medskip

\noindent
\textbf{Quotient structure.}
By definition,
\[
E^{\mathrm{proj}}(G)
=
E^{\mathrm{sph}}(G)/\mathbb Z_2
\]
is equipped with the quotient diffeology induced by the projection
\[
\pi:E^{\mathrm{sph}}(G)\to E^{\mathrm{proj}}(G).
\]

Since the action is smooth and free, the quotient map defines a principal
\(\mathbb Z_2\)-bundle in the diffeological sense.

\medskip

We now construct the characteristic class.

Choose a \(D\)-open cover
\[
(U_a)_a
\]
of \(E^{\mathrm{proj}}(G)\) over which the principal bundle admits smooth local
sections
\[
s_a:U_a\to E^{\mathrm{sph}}(G).
\]

On overlaps
\[
U_a\cap U_b,
\]
there exists a unique locally constant function
\[
g_{ab}:U_a\cap U_b\to\mathbb Z_2
\]
such that
\[
s_b=g_{ab}\cdot s_a.
\]

Since \(\mathbb Z_2\) is discrete, the family \((g_{ab})\) defines a
Čech \(1\)-cocycle with values in \(\mathbb Z_2\). Its cohomology class
\[
[(g_{ab})]
\in
\check H^1(E^{\mathrm{proj}}(G);\mathbb Z_2)
\]
depends only on the isomorphism class of the principal bundle.

Under the standard identification between principal \(\mathbb Z_2\)-bundles
and degree-one cohomology classes, this cocycle defines the characteristic
class
\[
\alpha_{\mathrm{sph}}
\in
H^1(E^{\mathrm{proj}}(G);\mathbb Z_2).
\]
\end{proof}

\end{document}
